% ============ §12 Krull rings (原书 pp.86-91) ============
\section{Krull rings}\label{sec:12}\index{Krull rings, section 12}

Let $A$ be an integral domain and $K$ its field of fractions. We write
$K^*$ for the multiplicative group of $K$. We say that $A$ is a \emph{Krull ring}\index{Krull ring} if there is
a family $\mathcal{F} = \{R_\lambda\}_{\lambda \in \Lambda}$ of DVRs of $K$ such that the following two conditions
hold, where we write $v_\lambda$ for the normalised additive valuation correspond-
ing to $R_\lambda$:
\begin{enumerate}[(1)]
  \item $A = \bigcap_\lambda R_\lambda$;
  \item for every $x \in K^*$ there are at most a finite number of $\lambda \in \Lambda$ such that
        $v_\lambda(x) \neq 0$.
\end{enumerate}
The family $\mathcal{F}$ of DVRs is said to be a \emph{defining family} of $A$. Since DVRs are
completely integrally closed (see Ex.~\ref{ex:9.5-ex}), so are Krull rings. If $A$ is a Krull
ring then for any subfield $K' \subset K$ the intersection $A \cap K'$ is again Krull.

\begin{theorem}\label{thm:12.1}
If $A$ is a Krull ring and $S \subset A$ a multiplicative set, then $A_S$ is
again Krull. If $\mathcal{F} = \{R_\lambda\}_{\lambda \in \Lambda}$ is a defining family of $A$ then the subfamily
$\{R_\lambda\}_{\lambda \in \Gamma}$, where $\Gamma = \{\lambda \in \Lambda \mid R_\lambda \supset A_S\}$, is a defining family of $A_S$.
\end{theorem}

\begin{proof}
Setting $m_\lambda$ for the maximal ideal of $R_\lambda$, we have
\[
  \bigcap_{\lambda \in \Gamma} m_\lambda = 0.
\]
Let $0 \neq x \in \bigcap_{\lambda \in \Gamma} R_\lambda$; there are at most finitely many $\lambda \in \Lambda$ such that $v_\lambda(x) < 0$;
let $\Delta = \{\lambda_1, \ldots, \lambda_n\}$ be the set of these. If $s \in S$ then $s \notin m_\lambda$, hence we can find
$t_\lambda \in m_\lambda \cap S$. Replacing $t_\lambda$ by a suitable power, we can assume that $v_\lambda(t_\lambda x) \ge 0$.
We then set $t = \prod_{\lambda \in \Delta} t_\lambda$, so that for every $\lambda \in \Lambda$ we have $v_\lambda(tx) \ge 0$, and
therefore $tx \in A$; but on the other hand $t \in S$ so that $x \in A_S$, and we have
proved that $A_S \supset \bigcap_{\lambda \in \Gamma} R_\lambda$. The opposite inclusion is obvious. The finiteness
condition (2) holds for $A$, so also for the subset $\Gamma$.
\end{proof}

Krull rings defined by a finite number of DVRs have a simple structure.

\begin{lemma}[Nagata]\label{lem:12.1}
Let $K$ be a field and $R_1, \ldots, R_n$ valuation rings of $K$;
set $A = \bigcap R_i$. Then for any given $a \in K$ there exists a natural number $s \ge 2$
such that
\[
  (1 + a + \cdots + a^{s-1})^{-1} \quad \text{and} \quad a \cdot (1 + a + \cdots + a^{s-1})^{-1}
\]
both belong to $A$.
\end{lemma}

\begin{proof}
We consider separately each $R_i$. Note first that $(1 - a)(1 + a + \cdots
+ a^{s-1}) = 1 - a^s$. If $a \notin R_i$ then $a^{-1} \in m_i$, and any $s \ge 2$ will do. If $a \in R_i$ then
provided that there does not exist $t$ such that $1 - a^t \in m_i$, any $s \ge 2$ will do.
If $1 - a \in m_i$ then any $s$ which is not a multiple of the characteristic of $R_i/m_i$
will do. If on the other hand $1 - a \notin m_i$ but $1 - a^t \in m_i$ for some $t > 2$, letting $t_0$
be the smallest value of $t$ for which this happens, we see that $1 - a^s \in m_i$
only for $s$ multiples of $t_0$, so that we only have to avoid these. Thus for each $i$
the bad values of $s$ (if any) are multiples of some number $d_i > 1$, so that
choosing $s$ not divisible by any of these $d_i$ we get the result.
\end{proof}

\begin{theorem}\label{thm:12.2}
Let $K$ be a field and $R_1, \ldots, R_n$ valuation rings of $K$ such
that $R_i \nsubseteq R_j$ for $i \neq j$; set $m_i = \rad(R_i)$. Then the intersection $A = \bigcap_{i=1}^{n} R_i$ is
a semilocal ring, having $p_i = m_i \cap A$ for $i = 1, \ldots, n$ as its only maximal
ideals; moreover $A_{p_i} = R_i$. If each $R_i$ is a DVR then $A$ is a principal ideal
ring.
\end{theorem}

\begin{proof}
The inclusion $A_{p_i} \subset R_i$ is obvious. For the opposite inclusion, let
$a \in R_i$; choosing $s \ge 2$ as in the lemma, and setting $u = (1 + a + \cdots + a^{s-1})^{-1}$
we get $u \in A$ and $au \in A$. Obviously $u$ is a unit of $R_i$, so that $u \in A - p_i$
and $a = (au)/u \in A_{p_i}$. This proves that $A_{p_i} = R_i$. It follows from this
that there are no inclusions among $p_1, \ldots, p_n$. If $I$ is an ideal of $A$ not
contained in any $p_i$ then (by Ex.~\ref{ex:1.6-ex}) there exists $x \in I$ not contained in
$\bigcup_{i=1}^{n} p_i$; then $x$ is a unit in each $R_i$, and hence in $A$, so that $I = A$. Thus
$p_1, \ldots, p_n$ are all the maximal ideals of $A$.

If each $R_i$ is a DVR then we have $m_i \neq m_i^{(2)}$, and hence $p_i \neq p_i^{(2)}$ (where $p^{(2)}$
denotes $p^2A_p \cap A$). Thus there exists $x_i \in p_i$ such that $x_i \notin p_i^{(2)}$, and $x_i \notin p_j$ for
$i \neq j$; then $p_i = x_iA$. If $I$ is any ideal of $A$ and $IR_i = x_i^{e_i}R_i$ for $i = 1, \ldots, n$ then it
is easy to see that $I = x_1^{e_1} \cdots x_n^{e_n}A$.
\end{proof}

If a Krull ring $A$ is defined by an infinite number of DVRs then the
defining family of DVRs is not necessarily unique, but the following
theorem tells us that among them there is a minimal family.

\begin{theorem}\label{thm:12.3}
Let $A$ be a Krull ring, $K$ its field of fractions, and $p$ a height $1$
prime ideal of $A$; then if $\mathcal{F} = \{R_\lambda\}_{\lambda \in \Lambda}$ is a family of DVRs of $K$ defining $A$, we
must have $A_p \in \mathcal{F}$. If we set $\mathcal{F}_0 = \{A_p \mid p \in \Spec A \text{ and } \htop p = 1\}$ then $\mathcal{F}_0$ is a
defining family of $A$. Thus $\mathcal{F}_0$ is the minimal defining family of DVRs of $A$.
\end{theorem}

\begin{proof}
By Theorem 1, $A_p$ is a Krull ring defined by the subfamily $\mathcal{F}_1 =
\{R_\lambda \mid A_p \subset R_\lambda\} \subset \mathcal{F}$; if $A_p \subset R_\lambda$ then the elements of $A - p$ are units of $R_\lambda$,
so that $p \supset m_\lambda \cap A$. If $m_\lambda \cap A = (0)$ then $R_\lambda \supset K$ which is a contradiction,
hence $m_\lambda \cap A \neq (0)$; since $\htop p = 1$, we must have $p = m_\lambda \cap A$. Thus if we fix
some $0 \neq x \in p$, then $v_\lambda(x) > 0$ for all $R_\lambda \in \mathcal{F}_1$, and hence $\mathcal{F}_1$ is a finite set,
Now by the previous theorem and Ex.~\ref{ex:10.5-ex}, the elements of $\mathcal{F}_1$ correspond
bijectively with maximal ideals of $A_p$, and $\mathcal{F}_1$ has just one element $A_p$. Thus
$A_p \in \mathcal{F}$, in other words $\mathcal{F}_0 \subset \mathcal{F}$.

To prove that $\mathcal{F}_0$ is a defining family of DVRs of $A$ it is enough to show that
$A \supset \bigcap_{\htop p = 1} A_p$. That is, it is enough to prove the implication
\[
  \text{for } a, b \in A \text{ with } a \neq 0, \quad b \in aA_p \text{ for all } A_p \in \mathcal{F}_0 \Rightarrow b \in aA.
\]
As one sees easily, this is equivalent to saying that $aA$ can be written as
the intersection of height $1$ primary ideals. The set of $R_\lambda \in \mathcal{F}$ such that
$aR_\lambda \neq R_\lambda$ is finite, so we write $R_1, \ldots, R_r$ for this. If we set
\[
  aR_i \cap A = q_i \quad \text{and} \quad \rad(R_i) \cap A = p_i
\]
then $q_i$ is a primary ideal belonging to $p_i$ for each $i$, and $aA = q_1 \cap \cdots \cap q_r$.
Eliminating redundant terms from this expression, we get an irredundant
expression, say $aA = q_1 \cap \cdots \cap q_r$. It is enough to show that then $\htop p_i = 1$
for $1 \le i \le r$. By contradiction, suppose that $\htop p_1 > 1$. By Theorem 1, $A_{p_1}$ is
a Krull ring with defining family $\mathcal{F}' = \{R_\lambda \in \mathcal{F} \mid A_{p_1} \subset R_\lambda\}$, but is not itself
a DVR, and hence by Theorem 2, $\mathcal{F}'$ is infinite. Thus there exists $R' \in \mathcal{F}'$
such that $aR' = R'$; we set $p' = \rad(R') \cap A$. We have $a \notin p'$, and $A_{p'} \subset R'$
implies that $p' \subset p_1$. Now by assumption $aA = q_1 \cap \cdots \cap q_r$ is irredundant, and $R_1$ is
a DVR, so that $(\rad(R_1))^\nu \subset aR_1$ for some $\nu > 0$, and hence $p_1^\nu \subset q_1$.
Therefore there exists an $i \ge 0$ such that
\[
  aA \neq p'^i \cap q_1 \cap \cdots \cap q_r \quad \text{and} \quad aA \supset p'^{i+1} \cap q_1 \cap \cdots \cap q_r.
\]
Hence there exist $b \in A$ such that $b \notin aA$ but $bp'^i \subset aA$. In particular $bp' \subset aA$,
but since $a$ is a unit of $R'$ we have
\[
  (b/a)p' \subset A \cap \rad(R') = p'.
\]
Taking $0 \neq c \in p'$ then for every $n > 0$ we have $(b/a)^nc \in p' \subset A$, and since $A$ is
completely integrally closed, $b/a \in A$. This is a contradiction, and it proves
that $\htop p_i = 1$ for $1 \le i \le r$.
\end{proof}

\begin{corollary*}
Let $A$ be a Krull ring and $\mathcal{P}$ the set of height $1$ prime ideals of $A$.
For $0 \neq a \in A$ set $v_p(a) = n_p$; then
\[
  aA = \bigcap_{p \in \mathcal{P}} p^{(n_p)},
\]
where $p^{(n)}$ denotes the symbolic $n$th power $p^nA_p \cap A$.
\end{corollary*}

\begin{proof}
According to the theorem we have $aA = \bigcap_p (aA_p \cap A)$,
but $aA_p = p^{n_p}A_p$, so that $aA_p \cap A = p^{(n_p)}$.
\end{proof}

\begin{theorem}\label{thm:12.4}
\begin{enumerate}[(i)]
  \item A Noetherian normal domain is a Krull ring.
  \item Let $A$ be an integral domain, $K$ its field of fractions, and $L$ an extension
        field of $K$. If $(A_i)_{i \in I}$ is a family of Krull rings contained in $L$ satisfying
        the two conditions (1) $A = \bigcap A_i$ and (2) given any $0 \neq a \in A$ we have
        $aA_i = A_i$ for all but finitely many $i$, then $A$ is a Krull ring.
  \item If $A$ is a Krull ring then so is $A[X]$ and $A[[X]]$.
\end{enumerate}
\end{theorem}

\begin{proof}
(i) This follows from Theorem~\ref{thm:11.5} and the fact that for any non-
zero $a \in A$ there are only finitely many height $1$ prime ideals containing $aA$
(because these are the prime divisors of $aA$).

(ii) is easy, and we leave it to the reader.

(iii) $K[X]$ is a principal ideal ring and therefore a Krull ring. Moreover,
if we let $\mathcal{P}$ be the set of height $1$ prime ideals of $A$ then for $p \in \mathcal{P}$
the ideal $p[X]$ is prime in $A[X]$, and by Theorem~\ref{thm:11.2}, (3), the local ring
$A[X]_{p[X]}$ is a DVR of $K(X)$. (If we write $v$ for the additive valuation of $K$
corresponding to the valuation ring $A_p$, we can extend $v$ to an additive
valuation of $K(X)$ by setting $v(F(X)) = \min\{v(a_i)\}$ for a polynomial
$F(X) = a_0 + a_1X + \cdots + a_rX^r$ (with $a_i \in K$),
and $v(F/G) = v(F) - v(G)$ for a rational function $F(X)/G(X)$; then the
valuation ring of $v$ in $K(X)$ is $A[X]_{p[X]}$.) Now we have
\[
  K[X] \cap A[X]_{p[X]} = A_p[X]
\]
(prove this!), and so
\[
  A[X] = \bigcap_{p \in \mathcal{P}} A_p[X],
\]
by (ii) this is a Krull ring.

Now for $A[[X]]$, let $\{R_\lambda\}_{\lambda \in \Lambda}$ be a family of DVRs of $K$ defining $A$; then
inside $K[[X]]$ we have $A[[X]] = \bigcap_\lambda R_\lambda[[X]]$; also by Ex.~\ref{ex:9.5-ex}, $R_\lambda[[X]]$ is an
integrally closed Noetherian ring, and is therefore a Krull ring by (i).
However, we cannot use (ii) as it stands, since $X$ is a non-unit of all the rings
$R_\lambda[[X]]$, so we set $R_\lambda[[X]][X^{-1}] = B_\lambda$ and note that $A[[X]] = K[[X]] \cap
(\bigcap_\lambda B_\lambda)$; now the hypothesis in (ii) is easily verified. Indeed,
\[
  F(X) = a_rX^r + a_{r+1}X^{r+1} + \cdots \in A[[X]] \quad \text{with } a_r \neq 0
\]
is a non-unit of $B_\lambda$ if and only if $a_r$ is a non-unit of $R_\lambda$, and there are only
finitely many such $\lambda$. Therefore $A[[X]]$ is a Krull ring.
\end{proof}

\begin{remark}
Note that the field of fractions of $A[[X]]$ is in general smaller
than the field of fractions of $K[[X]]$.
\end{remark}

\begin{remark}
The $B_\lambda$ occurring above are Euclidean rings ([B7], \S1, Ex. 9).
\end{remark}

\begin{theorem}\label{thm:12.5}
The notions of Dedekind ring and one-dimensional Krull
ring coincide.
\end{theorem}

\begin{proof}
A Dedekind ring is a normal Noetherian domain, and therefore a
Krull ring. Conversely, if $A$ is a one-dimensional Krull ring, let us prove that
$A$ is Noetherian. Let $I$ be a non-zero ideal of $A$, and let $0 \neq a \in I$. If we can
prove that $A/aA$ is Noetherian then $I/aA$ is finitely generated, and thus so is
$I$. By the corollary of Theorem 3 we can write $aA = q_1 \cap \cdots \cap q_r$, where
$q_i$ are symbolic powers of prime ideals $p_i$ and $p_i \neq p_j$ if $i \neq j$; now since
$\dim A = 1$ each $p_i$ is maximal and we have
\[
  A/aA = A/q_1 \times \cdots \times A/q_r
\]
by Theorem~\ref{thm:1.3} and Theorem~\ref{thm:1.4}. But $A/q_i$ is a local ring with maximal
ideal $p_i/q_i$, and hence $A/q_i \cong A_{p_i}/q_iA_{p_i}$; now since each $A_{p_i}$ is a DVR, $A/aA$ is
Noetherian (in fact even Artinian). Hence $A$ is one-dimensional Noetherian
integral domain, and is normal, and is therefore a Dedekind ring.
\end{proof}

\begin{theorem}\label{thm:12.6}
Let $A$ be a Krull ring, $K$ its field of fractions, and write $\mathcal{P}$
for the set of height $1$ prime ideals of $A$. Suppose given any $p_1, \ldots, p_r \in \mathcal{P}$
and $e_1, \ldots, e_r \in \Z$. Then there exists $x \in K$ satisfying
\[
  v_{p_i}(x) = e_i \quad \text{for } 1 \le i \le r
\]
and
\[
  v_p(x) \ge 0 \quad \text{for all } p \in \mathcal{P} - \{p_1, \ldots, p_r\}.
\]
Here $v_{p_i}$ and $v_p$ stand for the normalised additive valuations of $K$
corresponding to $p_i$ and $p$.
\end{theorem}

\begin{proof}
If $y_i \in A$ is chosen so that $y_i \in p_i$ but $y_i \notin p_i^{(2)}$ and $y_i \notin p_j$ for
$j \neq i$, then $v_{p_j}(y_i) = \delta_{ij}$ for $1 \le i, j \le r$. Similarly we choose $y_1, \ldots, y_r \in A$ such
that $v_{p_i}(y_j) = \delta_{ij}$. Then we set
\[
  y = \prod_{i=1}^{r} y_i^{e_i};
\]
let $p_1', \ldots, p_s'$ be all the primes $p \in \mathcal{P} - \{p_1, \ldots, p_r\}$ for which $v_p(y) \neq 0$.
Then choosing for each $j = 1, \ldots, s$ an element $t_j \in p_j'$ not belonging to
$p_1 \cup \cdots \cup p_r$, and taking $\nu$ to be sufficiently large, we see that
\[
  x = y(t_1 \cdots t_s)^\nu
\]
satisfies the requirements of the theorem.
\end{proof}

\begin{theorem}[Y. Mori and J. Nishimura]\label{thm:12.7}
Let $A$ and $\mathcal{P}$ be as in the previous
theorem. If $A/p$ is Noetherian for every $p \in \mathcal{P}$ then $A$ is Noetherian.
\end{theorem}

\begin{proof}[Proof (J. Nishimura)]
As in the proof of Theorem 5, it is enough to show
that $A/p^{(n)}$ is Noetherian for $p \in \mathcal{P}$ and any $n > 0$. If $n = 1$ this holds by
hypothesis. For $n > 1$ we proceed as follows. Applying the previous
theorem with $r = 1$ and $e = -1$ we can find an element $x$ in the field of
fractions of $A$ such that $v_p(x) = 1$ and $v_q(x) < 0$ for every other $q \in \mathcal{P}$. Set
$B = A[x]$. If $y \in p$ then $y/x \in A$ so that $y \in xB$, and conversely $B \subset A_p$, and
$xB \subset pA_p$, so that $p = xB \cap A$. Moreover, $B = A + xB$ and so
\[
  B/xB = A/p.
\]
Now $x^iB/x^{i+1}B \cong B/xB$ for each $i$, so that by induction on $i$ we see that
$B/x^iB$ is a Noetherian $B$-module for each $i$, and hence a Noetherian ring.
Now we have
\[
  A/(x^nB \cap A) \hookrightarrow B/x^nB,
\]
and $B/x^nB$ is finite over $A/(x^nB \cap A)$, being generated by the images of
$1, x, \ldots, x^{n-1}$, so that by the Eakin--Nagata theorem (Theorem~\ref{thm:3.7}),
$A/(x^nB \cap A)$ is a Noetherian ring; therefore its quotient $A/p^{(n)}$ is also
Noetherian.
\end{proof}

\begin{remark}
If $A$ is a Noetherian integral domain and $K$ its field of fractions,
then the integral closure of $A$ in $K$ is a possibly non-Noetherian Krull ring
([N1], (33.10)). This was proved by Y. Mori (1952) in the local case, and
in the general case by M. Nagata (1955). Theorem~\ref{thm:12.7} was proved by
Mori [1] in 1955 as a theorem on the integral closure of Noetherian
rings. His proof was correct (in spite of a number of easily rectifiable
inaccuracies), and was an extremely interesting piece of work, but due to
its difficulty, and the fact that it appeared in an inaccessible journal, the
result was practically forgotten. After Marot [1], (1973) applied it
successfully, Mori's work attracted attention once more, and J. Nishimura
[1], (1975) reformulated the result as above as a theorem on Krull rings
and gave an elegant proof.

More results on Krull rings can be found in [N1], [B7], [F], among
others.
\end{remark}

% ============ 习题 ============
\section*{Exercises to \S\ref{sec:12}}
\emph{Prove the following propositions.}

\begin{exercise}\label{ex:12.1-ex}
Let $K \subset L$ be a finite extension of fields, and $R$ a valuation ring of $K$. Then
there are a finite number of valuation rings of $L$ dominating $R$, and if $L$ is a
normal extension of $K$ then these are all conjugate to one another under
elements of the Galois group $\Aut_K(L)$.
\end{exercise}

\begin{exercise}\label{ex:12.2-ex}
Let $R$ be a valuation ring of a field $K$, and let $K \subset L$ be a (possibly infinite)
algebraic extension; write $\bar{R}$ for the integral closure of $R$ in $L$. Then the
localisation of $\bar{R}$ at a maximal ideal is a valuation ring dominating $R$, and
conversely every valuation ring of $L$ dominating $R$ is obtained in this way.
\end{exercise}

\begin{exercise}\label{ex:12.3-ex}
Let $A$ be a Krull ring, $K$ its field of fractions, and $K \subset L$ a finite extension
field; if $B$ is the integral closure of $A$ in $L$, then $B$ is also a Krull ring.
\end{exercise}

\begin{exercise}\label{ex:12.4-ex}
Let $A$ be an integral domain and $K$ its field of fractions. For $I$ a fractional
ideal, write $\tilde{I} = (I^{-1})^{-1}$. If $I = \tilde{I}$ we say that $I$ is \emph{divisorial}. If $A$ is a Krull
ring, then an ideal of $A$ is divisorial if and only if it can be expressed as the
intersection of a finite number of height $1$ primary ideals.
\end{exercise}
